(1) 平行太阳光射入折射率 \(1<n<2\) 的球形水滴，入射角 \(i\in(0,\pi/2)\)，折射角 \(r\) 满足 \(\sin i=n\sin r\)。光线经两次折射、一次内反射离开水滴。求相对初始光线的总偏向角 \(D(i)\)。
(2) 推导产生光线聚集的偏向角驻点条件，求其对应 \(\sin^2 i_*\)。
(3) 定义从反太阳方向测量的主虹角半径 \(\alpha=\pi-D(i_*)\)，写出 \(\alpha\) 随折射率 \(n\) 的关系。
(4) 用普通色散 \(dn/d\lambda<0\) 解释红光在外、紫光在内的次序，并说明几何光学预言的强度奇性为何被波动光学平滑。
